\documentclass[11pt]{article}

\usepackage{fontspec}
\usepackage{polyglossia}
\setmainlanguage{french}
\setmainfont{Latin Modern Roman}

\usepackage[margin=2.6cm]{geometry}
\usepackage{booktabs}
\usepackage{amsmath}
\usepackage{amssymb}
\usepackage{enumitem}
\usepackage{titlesec}
\usepackage[hidelinks]{hyperref}
\usepackage{tcolorbox}
\tcbuselibrary{breakable}

\titleformat{\section}{\Large\bfseries}{\thesection}{0.7em}{}
\titleformat{\subsection}{\large\bfseries}{\thesubsection}{0.6em}{}

\newtcolorbox{remarque}[1][Remarque]{breakable, colback=black!3, colframe=black!40,
  boxrule=0.4pt, arc=1pt, title=#1, fonttitle=\bfseries, coltitle=black}
\newtcolorbox{cle}[1][À retenir]{breakable, colback=blue!4, colframe=blue!40!black,
  boxrule=0.5pt, arc=1pt, title=#1, fonttitle=\bfseries, coltitle=blue!40!black}
\newtcolorbox{question}[1]{breakable, colback=black!2, colframe=black!55,
  boxrule=0.4pt, arc=1pt, title=Question probable~: #1, fonttitle=\bfseries\itshape,
  coltitle=black}

\setlength{\parskip}{0.5em}
\setlength{\parindent}{0pt}

\title{Fiche de maîtrise, Résultat A\\[0.3em]
\large Pour l'expliquer sans notes, pas pour le relire}
\author{Tidiane Diallo}
\date{5 novembre 2026}

\begin{document}
\maketitle

\begin{remarque}[Objectif de cette fiche]
Ce n'est pas un résumé du document envoyé au directeur
(\texttt{RESULTAT\_A\_COMPLET\_05\_NOVEMBRE\_2026.tex}). C'est le fil
narratif qui doit rester en tête pour pouvoir expliquer le résultat A à
l'oral, reconstruire chaque étape sans notes, et répondre aux questions les
plus probables sans hésiter.
\end{remarque}

\section{L'idée en une phrase}

Carlini et al. savent tout récupérer d'un réseau ReLU en hard-label sauf la
dernière couche, parce que cette couche n'a pas d'activation et donc pas de
coude à détecter. Mon article de départ (Canales-Martínez et Santos) comble
ce trou, mais seulement si le reste du réseau est en ReLU. Le résultat A
montre que cette restriction est inutile : la méthode qui récupère la
dernière couche ne regarde jamais l'activation cachée, elle regarde
uniquement la règle de décision finale, un argmax, qui est la même quelle
que soit l'activation.

\begin{cle}
Si on me coupe au milieu d'une explication, c'est la phrase à ressortir en
premier : l'indépendance à l'activation n'est pas un espoir, c'est une
conséquence directe de trois hypothèses qui ne mentionnent jamais
l'activation.
\end{cle}

\section{Les trois hypothèses, en profondeur}

\textbf{(H1)} Au point de transition $p$, $A_i \cdot y + b_i = A_j \cdot y +
b_j$. C'est juste la définition d'une frontière de décision entre deux
classes : le point où les deux scores sont à égalité parfaite.

\textbf{(H2)} $y$ est calculable. C'est l'hypothèse qui fait le pont avec
EUROCRYPT 2025 : eux savent extraire les couches cachées en ReLU hard-label,
donc à partir de leur résultat, l'attaquant peut calculer $y = f_{1..r}(x)$
pour n'importe quel $x$ qu'il choisit, même sans connaître l'activation en
interne, puisqu'il a déjà les poids. Cette hypothèse, c'est « je prends le
relais là où l'article précédent s'arrête ». Si on demande pourquoi je
suppose ça~: ce n'est pas une triche, c'est la continuation logique d'une
chaîne d'attaque déjà publiée.

\textbf{(H3)} C'est la clé de tout. L'oracle ne renvoie que $\arg\max_k$.
Deux propriétés en découlent.

\begin{cle}[L'exemple à retenir par cœur]
Scores $= (1, 5, 3)$. $\arg\max = $ indice 1 (valeur 5).
\begin{itemize}
  \item Ajouter $10$ à tous : $(11, 15, 13)$. $\arg\max$ toujours indice 1.
    \textbf{Translation inoffensive.}
  \item Multiplier tous par $2$ : $(2, 10, 6)$. $\arg\max$ toujours indice 1.
    \textbf{Échelle positive inoffensive.}
  \item Multiplier tous par $-1$ : $(-1, -5, -3)$. $\arg\max$ devient
    l'indice 0 (valeur $-1$), qui était l'indice du \textbf{minimum}
    original. \textbf{Échelle négative : l'argmax devient un argmin.}
\end{itemize}
Ces trois lignes expliquent, à elles seules, pourquoi « accord avant
correction = $0/5000$ » apparaît dans le tableau de résultats pour ReLU,
GELU et SiLU.
\end{cle}

\section{Pourquoi ça prouve l'indépendance à l'activation}

Aucune des trois hypothèses ne parle de $\sigma$, l'activation. H1 est une
égalité de scores. H2 dit juste que $y$ est calculable, peu importe
comment. H3 porte sur l'argmax, qui s'applique après la dernière couche,
jamais avant. Une fois que $y$ est traité comme une boîte noire calculable,
le problème qui reste, retrouver un classifieur linéaire à partir de son
argmax, est un problème purement algébrique qui ne voit plus ReLU, GELU,
SiLU ou tanh.

\begin{question}{« Ce n'est pas juste une évidence une fois qu'on lit les hypothèses ? »}
Si, en partie, et c'est assumé dans le document. Mais personne ne l'avait
vérifié explicitement sur quatre activations avec des mesures réelles avant
ce travail, et la correction de l'ambiguïté de signe en sortie
multi-classes va au-delà de ce que l'article de départ traite (lui ne
discute que la sortie binaire). Ça, c'est une extension, pas seulement une
vérification.
\end{question}

\section{Compter les degrés de liberté}

$36$ coefficients à retrouver ($4$ classes $\times$ ($8$ poids $+$ $1$
biais)). Mais l'attaquant ne peut jamais distinguer deux couches de sortie
qui, après translation ou mise à l'échelle positive, donnent le même argmax
partout. C'est la même limite que « je ne peux pas savoir si tu mesures en
degrés décalés d'une constante, seule la différence de température
compte ».

\begin{itemize}
  \item $9$ degrés invisibles pour la translation (décaler les 8 colonnes
    de $A$ et le biais ne change rien à l'écart entre deux scores).
  \item $1$ degré invisible pour l'échelle globale.
\end{itemize}

$36 - 10 = 26$.

\begin{cle}
« $26/26$ » dans le tableau de résultats n'est pas une bonne performance
par chance, c'est la preuve que la méthode atteint exactement le plafond
mathématiquement atteignable, ni plus ni moins. C'est un argument de
correction, pas de qualité.
\end{cle}

\section{Le fil narratif des dix étapes}

Pas les chiffres, juste la logique de chaque étape, pour pouvoir la
reconstruire sans notes.

\begin{enumerate}[leftmargin=1.4em]
  \item \textbf{Construire le réseau cible.} On connaît sa vraie couche de
    sortie pour pouvoir vérifier après coup. L'attaquant, lui, ne la voit
    jamais.
  \item \textbf{Trouver un point de transition par dichotomie.} Deux points
    donnant des classes différentes, on coupe en deux, on regarde de quel
    côté est la frontière, on recommence. Chaque itération divise
    l'incertitude par 2. Après une cinquantaine d'itérations, on tombe sur
    le mur de la précision machine ($2^{-52}$) : continuer ne sert à rien,
    le flottant ne peut plus être plus précis.
  \item \textbf{Identifier les deux classes en jeu} en sondant juste à
    gauche et à droite du point trouvé.
  \item \textbf{Construire une équation linéaire} à partir de ce point :
    $(A_i - A_j)\cdot y + (b_i - b_j) = 0$. Un point de transition donne
    exactement une ligne du système. C'est le cœur mécanique de toute la
    méthode.
  \item \textbf{Répéter jusqu'à avoir assez de points} (156 ici, soit
    $6 \times 26$, une marge de sécurité, pas le strict minimum).
  \item \textbf{Vérifier le rang avant de résoudre.} La discipline la plus
    importante de tout le protocole : sans cette vérification,
    \texttt{lstsq} renvoie quand même une solution, mais une solution
    fausse, de norme minimale, pas la bonne, sans prévenir. Vérifier que le
    noyau mesuré correspond exactement aux générateurs attendus
    (translations + échelle) transforme « ça a tourné sans erreur » en
    « c'est mathématiquement correct ».
  \item \textbf{Résoudre en fixant arbitrairement 10 variables.} Un choix
    de représentant dans la classe d'équivalence, comme choisir où mettre
    le zéro d'un thermomètre.
  \item \textbf{Corriger l'ambiguïté de signe.} Le moment le plus
    spectaculaire : avant correction, $0\,\%$ d'accord, bien pire que le
    hasard ($25\,\%$). Ça fait peur si on ne sait pas pourquoi, mais c'est
    exactement la signature attendue d'une échelle négative (l'exemple à
    trois nombres de la section~2). Une seule requête suffit à la
    détecter.
  \item \textbf{Valider sur deux distributions différentes}, pas une
    seule, pour être sûr que l'équivalence est fonctionnelle partout, pas
    seulement sur les données d'entraînement.
  \item \textbf{Regarder où part le coût} : $100\,\%$ dans la recherche des
    points de transition. C'est l'information la plus utile pour la suite
    de la thèse côté défense : c'est cette étape-là qu'il faut rendre
    difficile, pas la résolution du système, qui est instantanée.
\end{enumerate}

\section{Questions probables, réponses courtes}

\begin{question}{« Pourquoi le facteur d'échelle est négatif 3 fois sur 4 ? »}
Pur hasard du tirage aléatoire de $A_{1,1,\text{vrai}}$, rien de
systématique. Avec une graine différente, le mélange de signes
positifs/négatifs serait différent.
\end{question}

\begin{question}{« C'est quoi une classe d'équivalence ici, concrètement ? »}
L'ensemble de tous les $(A,b)$ qui donnent exactement le même argmax pour
toute entrée. On ne récupère jamais le $(A,b)$ exact, c'est mathématiquement
impossible sans information supplémentaire, on récupère un représentant de
sa classe, ce qui est tout ce que l'oracle peut révéler de toute façon.
\end{question}

\begin{question}{« Pourquoi GELU prend deux fois plus de temps ? »}
Coût de simulation de l'oracle (la fonction $\mathrm{erf}$ réimplémentée en
Python), pas un coût propre à l'attaque. Le nombre de requêtes, la vraie
mesure qui compte, est quasi identique entre les quatre activations.
\end{question}

\begin{question}{« Est-ce que ça marcherait sur un vrai réseau, pas un réseau jouet ? »}
C'est précisément la limite assumée dans le document : hypothèse H2 non
vérifiée de bout en bout, une seule graine testée, pas encore rejoué sur le
réseau CIFAR-10 de l'article de départ. Prochaine étape déjà identifiée.
\end{question}

\end{document}
